\documentclass[11pt]{letter}
\usepackage[margin=2.5cm]{geometry}
\usepackage{microtype}
\signature{A. Urias}
\address{Independent researcher\\C.P. 83220, Hermosillo, Sonora, Mexico\\alex@aurtech.mx}
\begin{document}
\begin{letter}{The Editors\\Forensic Science International: Digital Investigation}
\opening{Dear Editors,}

We submit \emph{Evidence-Monotone Audio Representation: Preventing Provenance
Promotion Across Derived Audio Transformations} for consideration as a research
article.

Current provenance standards can cryptographically bind assertions and
transformation history to an asset. The missing layer addressed here is whether
the evidence claim emitted by a derived representation is conservative with
respect to every source interval that representation actually contains. A clip
assembled from captured and generated material, then trimmed, resampled and
re-encoded, can emerge described as captured-only, with a valid signature, an
intact chain and a bit-identical waveform, because a derived asset inherits
from its parent and a minimal implementation summarises that parent from the
boundaries of the retained region. For an examiner reading such an
asset, the difference between that claim and the truth is the whole point.

We contribute a formal operator contract for derived media, a temporal audio
instantiation with declared kernel footprints and a footprint-monotonicity
result that makes the rule implementable against real codecs, a signal-transparent
reference implementation that requires no model training, a signed transport
built entirely inside C2PA's own extension points, and a frozen public
reproducibility artifact with one-command reproduction.

Four things are worth flagging to the editors directly, because we would rather
state them than have a reviewer find them.

First, we do not claim the underlying algebra as new. The non-amplification
property is an instance of a meet in a lattice, familiar from provenance
semirings and from information-flow models, and we say so in the Introduction,
the Formal model and the Discussion. The contribution is the executable temporal
contract, the transport and the measurements.

Second, the comparison policy is a constructed reference policy motivated by the
specification's own \texttt{parentOf} definition. It is not a description of any
shipping product, no product is named or evaluated, and the fixture rates are
reported as a mechanism stress rate rather than as a prevalence estimate.

Third, this is not a deepfake-detection paper. No model is trained, no acoustic
feature is used, and an unsigned or provenance-stripped asset is classified as
unverified rather than synthetic. A cryptographically valid signature proves
that an assertion was signed, not that it is true, and we state that limitation
in the abstract, the threat model, the discussion and the conclusion.

Fourth, the paper reports an independent reproduction that partly failed, and
reports it as a result. A colleague at another institution, who is acknowledged
but is not an author, ran the packaged release twice on unrelated hardware and a
different FFmpeg major version. Every headline count returned unchanged and
eight of the ten compared result files showed no deterministic difference at
all. Two declared kernel footprints turned out not to be adequate on his build,
one of them under-sized. We have not widened those declarations to make the
comparison pass; Section~7.11 reports what he measured, and his unmodified
outputs are released with the artifact.

The work fits the journal's stated primary pillar of digital evidence and
multimedia with the core qualities of provenance, integrity and authenticity,
and its ``scientific practices'' track: it is an attempt to make one specific
property of derived digital evidence falsifiable and reproducible rather than
assumed.

The manuscript is original, is not under consideration elsewhere, and all
authors have approved the submission. A declaration of competing interests and a
declaration of generative-AI use are included. All code, fixtures, corpora
build scripts and machine-readable results are public under the MIT licence,
with third-party data licences verified from primary sources.

\closing{Yours sincerely,}
\end{letter}
\end{document}
